\documentclass[11pt]{article}
\usepackage[margin=1in]{geometry}
\usepackage{amsmath,amssymb,amsthm}
\usepackage{graphicx}
\usepackage{hyperref}
\usepackage{xcolor}
\newtheorem{theorem}{Theorem}
\newtheorem{corollary}[theorem]{Corollary}
\newtheorem{lemma}[theorem]{Lemma}
\newcommand{\Lzero}{L^0(\mathcal F)}
\newcommand{\Tep}{\mathcal T_{\varepsilon,\lambda}}
\newcommand{\epi}{\operatorname{epi}}
\title{Closed Random Sublevels Force Epigraph Closure in Two Broad Regimes\\
\large A partial answer to Remark 2.14 of \emph{On random convex analysis}}
\author{Run \texttt{fa\_banach\_001}, lane 3 (GPT5.6)}
\date{9 August 2026}

\begin{document}
\maketitle

\begin{center}
\fcolorbox{black}{gray!10}{\parbox{0.91\linewidth}{
\textbf{Status.} Candidate partial result, likely valid, pending human review.
The full arbitrary-net question remains open here.}}
\end{center}

\section{Source question}

Guo, Zhang, Wu, Yang, Yuan, and Zeng ask the following in Remark 2.14 of
\cite{GuoEtAl}, PDF page 10.  Let $(E,\mathcal P)$ be a random locally convex
(RLC) module over a sigma-finite measure space $(\Omega,\mathcal F,\mu)$, and
let $f:E\to\overline{L}^0(\mathcal F)$ be proper and local.  Compare
\begin{align*}
(1)'\quad &\{x\in E:f(x)\leq r\}\text{ is $\Tep$-closed for every }r\in\Lzero,\\
(2)'\quad &\epi(f)=\{(x,r):f(x)\leq r\}\text{ is closed in }
(E,\Tep)\times(\Lzero,\Tep).
\end{align*}
The authors observe $(2)'\Rightarrow(1)'$ and state that they do not know
whether $(1)'\Rightarrow(2)'$ in general. Their Theorem 2.13 gives the
implication when the module and seminorm family have the countable
concatenation property.

\begin{figure}[ht]
\centering
\includegraphics[width=\textwidth]{figures/open_problem_crop.png}
\caption{Remark 2.14 in the source paper (PDF page 10).}
\end{figure}

\section{Idea of the proof}

The obstruction posed by random thresholds disappears for sequences.  If
$r_n\to r$ in probability, a subsequence converges almost everywhere.  Its
pointwise supremum is consequently a finite random variable.  On any
positive-measure set where $f(x)$ lies uniformly above $r$, one can discard a
finite initial segment so that this supremum remains strictly below $f(x)$ on
a positive subset.  All corresponding $x_n$ then lie in one fixed random
sublevel, whose sequential closedness gives the contradiction.

This mechanism yields sequential epigraph closure on every RLC module.  When
the module topology is first countable, sequential closure is ordinary
closure, proving the requested implication without countable concatenation.

There is a second mechanism for arbitrary closure nets.  If the module $E$
has the countable concatenation property, closedness of one fixed sublevel
provides a separating neighborhood involving only finitely many seminorms.
One selects countably many epigraph points with summable errors for those
seminorms, partitions the base by the sets on which their thresholds lie
below the fixed level, and pastes the points along that partition.  This proves
the implication without any countable-concatenation hypothesis on the
seminorm family $\mathcal P$.

\section{The sequential theorem}

We use the standard fact that the $\Tep$ topology on $\Lzero$ is metrizable
and is convergence in probability for an equivalent probability measure
$\widehat\mu$ (equivalently, convergence locally in measure for the original
sigma-finite measure), as recalled in \cite[Section 2]{GuoEtAl}.

\begin{lemma}[random upper envelope]
\label{lem:envelope}
Let $r_n,r\in\Lzero$ and $r_n\to r$ in $\Tep$.  Let $B\in\mathcal F$ have
positive measure and let $\delta>0$.  There are a subsequence
$(r_{n_j})$, an integer $N$, a positive-measure set $D\subset B$, and
$q\in\Lzero$ such that
\[
r_{n_j}\leq q\quad(j\geq N),
\qquad q\leq r+\delta\quad\text{on }D.
\]
\end{lemma}

\begin{proof}
Pass to a subsequence, still denoted $(r_{n_j})$, which converges to $r$
$\widehat\mu$-almost everywhere.  For $N\geq1$, put
\[
D_N=B\cap\bigcap_{j\geq N}\{r_{n_j}\leq r+\delta\}.
\]
Modulo a null set, $D_N\uparrow B$, so $D_N$ has positive measure for some
$N$.  After fixing such an $N$, set
\[
q=\sup_{j\geq N}r_{n_j}.
\]
The supremum is measurable and finite almost everywhere because the
underlying sequence converges almost everywhere.  Thus $q\in\Lzero$; it
dominates every retained $r_{n_j}$, and $q\leq r+\delta$ on $D_N$.
\end{proof}

\begin{theorem}[closed sublevels imply sequentially closed epigraph]
\label{thm:sequential}
Let $X$ be any topological space and
$f:X\to\overline L^0(\mathcal F)$.  Suppose that for every $q\in\Lzero$ the
sublevel
\[
C_q=\{x\in X:f(x)\leq q\}
\]
is sequentially closed in $X$.  Then $\epi(f)$ is sequentially closed in
$X\times(\Lzero,\Tep)$.
\end{theorem}

\begin{proof}
Let $(x_n,r_n)\in\epi(f)$ and suppose $(x_n,r_n)\to(x,r)$.  Thus
$f(x_n)\leq r_n$ for all $n$.  Assume for a contradiction that
$f(x)\nleq r$.  Then $\{f(x)>r\}$ has positive measure.  By taking a
countable union of uniform-gap sets, choose $\delta>0$ and a
positive-measure set $B$ such that
\[
f(x)>r+2\delta\quad\text{on }B.
\]
Apply Lemma~\ref{lem:envelope}.  For the resulting subsequence, tail, set
$D$, and random variable $q$, we have
\[
f(x_{n_j})\leq r_{n_j}\leq q\qquad(j\geq N).
\]
Hence $x_{n_j}\in C_q$ for all $j\geq N$.  This subsequence still converges
to $x$, and $C_q$ is sequentially closed, so $f(x)\leq q$.  On $D$, however,
\[
q\leq r+\delta<f(x),
\]
a contradiction on a set of positive measure.  Therefore $f(x)\leq r$ and
$(x,r)\in\epi(f)$.
\end{proof}

Theorem~\ref{thm:sequential} is stronger than needed in two respects: it
requires neither an $L^0$-module structure nor properness or locality of $f$.

\section{The countable-concatenation theorem}

\begin{theorem}[the module CCP alone is sufficient]
\label{thm:ccp}
Let $(E,\mathcal P)$ be an RLC module and suppose that $E$ has the countable
concatenation property; no such assumption is made on $\mathcal P$.  Let
$f:E\to\overline L^0(\mathcal F)$ be proper and local.  If every sublevel
$C_q=\{f\leq q\}$, $q\in\Lzero$, is $\Tep$-closed, then $\epi(f)$ is closed
in $(E,\Tep)\times(\Lzero,\Tep)$.
\end{theorem}

\begin{proof}
Write $\widehat\mu$ for the probability measure defining $\Tep$.  Suppose
that $(x,r)$ belongs to the closure of $\epi(f)$ but not to the epigraph.
There are a set $B\in\mathcal F$ of positive measure and a deterministic
$\delta>0$ such that
\[
f(x)>r+2\delta\quad\text{on }B.
\]
Choose $u\in\operatorname{dom}(f)$ and localize outside $B$ by putting
\[
\bar x=I_Bx+I_{B^c}u,
\qquad
\bar r=I_Br+I_{B^c}f(u),
\qquad
q=I_B(r+\delta)+I_{B^c}f(u).
\]
Locality shows that localization sends epigraph points to epigraph points,
and multiplication by indicators is continuous.  Hence
$(\bar x,\bar r)$ is still in the closure of $\epi(f)$.  Moreover
$f(\bar x)>q$ on $B$, so $\bar x\notin C_q$.

Since $C_q$ is closed, choose a finite nonempty $Q\subset\mathcal P$ and
$\varepsilon,\lambda>0$ such that
\begin{equation}
\label{eq:separating-neighborhood}
\bar x+N_0(Q,\varepsilon,\lambda)\quad\text{is disjoint from }C_q,
\end{equation}
where
\[
N_0(Q,\varepsilon,\lambda)
=\{v:\widehat\mu(\{\|v\|_Q<\varepsilon\})>1-\lambda\}.
\]
Choose positive numbers $\eta_n$ with $\eta_n\downarrow0$ and
$\sum_n\eta_n<\lambda$.  Using the closure of the epigraph and then applying
the preceding localization, choose $(z_n,t_n)\in\epi(f)$ such that
\begin{align}
\widehat\mu\{\|z_n-\bar x\|_Q\geq\varepsilon\}&<\eta_n,
\label{eq:space-errors}\\
\widehat\mu\{|t_n-\bar r|\geq\delta\}&<\eta_n,
\label{eq:level-errors}
\end{align}
and such that $z_n=u$ and $t_n=f(u)$ on $B^c$.  Let
$G_n=\{t_n\leq q\}$.  The equality on $B^c$ and
\eqref{eq:level-errors} give
\[
\widehat\mu(G_n^c)<\eta_n.
\]
Consequently $\bigcup_nG_n=\Omega$ modulo null sets, since the measure of
$\bigcap_nG_n^c$ is at most $\inf_n\eta_n=0$.  Define the measurable
partition
\[
D_1=G_1,
\qquad
D_n=G_n\setminus\bigcup_{k<n}G_k\quad(n\geq2).
\]

The CCP of $E$ supplies
\[
y=\sum_{n=1}^{\infty}I_{D_n}z_n\in E.
\]
For each $n$, locality gives
$I_{D_n}f(y)=I_{D_n}f(z_n)$.  Because $D_n\subset G_n$ and
$f(z_n)\leq t_n$, it follows that $f(y)\leq q$; hence $y\in C_q$.
On the other hand, locality of every random seminorm gives
\[
\{\|y-\bar x\|_Q\geq\varepsilon\}
=\bigcup_n\left(D_n\cap
\{\|z_n-\bar x\|_Q\geq\varepsilon\}\right)
\quad\text{modulo null sets}.
\]
By \eqref{eq:space-errors}, its probability is at most
$\sum_n\eta_n<\lambda$.  Thus
$y\in\bar x+N_0(Q,\varepsilon,\lambda)$, contradicting
\eqref{eq:separating-neighborhood}.
\end{proof}

Source Theorem 2.13 assumes that both $E$ and $\mathcal P$ have the
countable concatenation property.  Theorem~\ref{thm:ccp} removes the
assumption on $\mathcal P$ for the implication $(1)'\Rightarrow(2)'$.

\section{Affirmative subcases of Remark 2.14}

\begin{corollary}[first-countable case]
\label{cor:firstcountable}
Let $(E,\mathcal P)$ be an RLC module such that $(E,\Tep)$ is first
countable.  For any map $f:E\to\overline L^0(\mathcal F)$, if every random
sublevel $\{f\leq q\}$ is $\Tep$-closed, then $\epi(f)$ is closed in
$(E,\Tep)\times(\Lzero,\Tep)$.
\end{corollary}

\begin{proof}
Closed sublevels are sequentially closed, so
Theorem~\ref{thm:sequential} makes the epigraph sequentially closed.  The
factor $\Lzero$ is metrizable, and the product of a first-countable space
with a metrizable space is first countable.  In a first-countable space,
sequentially closed sets are closed.
\end{proof}

\begin{corollary}[countably generated RLC modules]
The implication $(1)'\Rightarrow(2)'$ in Remark 2.14 holds, without any
countable concatenation assumption, if the $\Tep$ topology of $E$ is
generated by a countable family of random seminorms.  In particular it holds
for every random normed module.
\end{corollary}

\begin{proof}
For a countable seminorm family, finite subfamilies together with positive
rational values of $\varepsilon$ and $\lambda$ give a countable local base
for $\Tep$.  Thus Corollary~\ref{cor:firstcountable} applies.  A random normed
module has a singleton defining seminorm family.
\end{proof}

\section{Why the full net problem remains}

When $E$ lacks the CCP, the proof of Theorem~\ref{thm:ccp} produces only a
formal countable paste, which need not be an element of $E$.  The sequential
proof also cannot simply replace ``subsequence'' by ``subnet.''  A net
may converge in probability while carrying arbitrarily high values on
measurable sets whose probabilities shrink and whose locations move with the
index.  A whole tail then need not possess a finite pointwise random upper
bound.  Moreover, for an uncountable seminorm family the $\Tep$ topology of
$E$ need not be first countable or sequential, so a point in the closure of
the epigraph need not be witnessed by an $E$-convergent sequence.

Thus the theorems settle the sequential content of the source question, all
first-countable cases, and all modules $E$ with the CCP (regardless of
$\mathcal P$).  They neither prove nor disprove the implication in the
remaining regime where $E$ lacks the CCP and its topology is genuinely
nonsequential.  This is why the packet remains classified as a partial result.

\section{Verification and novelty status}

The proofs use only: (i) the elementary almost-everywhere subsequence
principle for convergence in probability; (ii) finiteness of the pointwise
supremum of an almost-everywhere convergent real-valued sequence; and (iii)
the first-countable characterization of closed sets by sequences; and (iv)
locality plus summable probability errors for the countable-paste argument.
No computational or unproved dependency is present.

The run's lightweight indexes and local arXiv corpus were searched for the
paper id, the exact Remark 2.14 question, and close epigraph/sublevel and
countable-concatenation terminology.  Bounded web searches on 9 August 2026
found the source and background papers but no later resolution of either the
stated subcase or the full question.  This is evidence of plausible novelty,
not a comprehensive certification.

\paragraph{Human-review focus.}
Verify Lemma~\ref{lem:envelope}, especially the positive-measure tail set and
the almost-everywhere finiteness of $q$; verify the localization and good-set
partition in Theorem~\ref{thm:ccp}; check the summable seminorm-error estimate
for its pasted point; and verify that the chosen seminorm hypothesis indeed
supplies a countable $\Tep$ local base.

\begin{thebibliography}{9}
\bibitem{GuoEtAl}
T. Guo, E. Zhang, M. Wu, B. Yang, G. Yuan, and X. Zeng,
\emph{On random convex analysis}, arXiv:1603.07074 (2016), especially
Section 2 and Remark 2.14.
\end{thebibliography}

\end{document}
